\subsection{\texorpdfstring{模的计算 \(\operatorname{Tor}^{\mathbf{A}}(\mathbf{P},\mathbf{M})\) 与 \(\operatorname{Ext}_{\mathbf{A}}(\mathbf{M},\mathbf{N})\)}{模的计算 \textbackslash operatorname\{Tor\}\^{}\{\textbackslash mathbf\{A\}\}(\textbackslash mathbf\{P\},\textbackslash mathbf\{M\}) 与 \textbackslash operatorname\{Ext\}\_\{\textbackslash mathbf\{A\}\}(\textbackslash mathbf\{M\},\textbackslash mathbf\{N\})}}

设 M、N 为左 A-模，P 为右 A-模。此外，设 \(a: R \rightarrow M\) 是 M 的一个左分解， \(b: S \rightarrow P\) 为 P 的一个左分解，且 \(c: N \rightarrow E\) 为 N 的一个右分解。

由 X，第 49 页，命题 3 及 3 bis，存在复形态射 \(\alpha: L(M) \to R\) , \(\beta: L(P) \to S\) , \(\gamma: E \to I(N)\) 使得 \(a \circ \alpha = p_{M}\) , \(b \circ \beta = p_{P}\) , \(\gamma \circ c = e_{N}\) ；且 \(\alpha\) , \(\beta\) , \(\gamma\) 的同伦类仅依赖于给定的分解。由 X，第 64 页，命题 3 及 X，第 84 页，命题 3，态射

\[
\beta \otimes \alpha : L (P) \otimes_ {A} L (M) \rightarrow S \otimes_ {A} R,
\]

\[
\operatorname{Homgr} _ {\mathbf {A}} (\alpha , \gamma): \operatorname{Homgr} _ {\mathbf {A}} (\mathrm{R}, \mathrm{E}) \rightarrow \operatorname{Homgr} _ {\mathbf {A}} (\mathrm{L} (\mathrm{M}), \mathrm{I} (\mathrm{N}))
\]

仅依赖于给定的分解，从而通过取同调，得到次数为 0 的分次 k-同态

\[
\psi (\mathrm{S}, \mathrm{R}): \operatorname{Tor} ^ {\mathrm{A}} (\mathrm{P}, \mathrm{M}) \rightarrow \mathrm{H} (\mathrm{S} \otimes_ {\mathrm{A}} \mathrm{R}),
\]

\[
\varphi (\mathrm{R}, \mathrm{E}): \mathrm{H} \left(\operatorname{Homgr} _ {\mathrm{A}} (\mathrm{R}, \mathrm{E})\right)\rightarrow \operatorname{Ext} _ {\mathrm{A}} (\mathrm{M}, \mathrm{N}),
\]

与 α、β、γ 的选择无关。

例如，取 \(a\) , \(b\) , \(c\) 分别为 M、P、N 的恒等映射，便得到同态 \(\psi(\mathbf{P}, \mathbf{M}) : \operatorname{Tor}^{\mathbf{A}}(\mathbf{P}, \mathbf{M}) \to \mathbf{P} \otimes_{\mathbf{A}} \mathbf{M}\) 与 \(\varphi(\mathbf{M}, \mathbf{N}) : \operatorname{Hom}_{\mathbf{A}}(\mathbf{M}, \mathbf{N}) \to \operatorname{Ext}_{\mathbf{A}}(\mathbf{M}, \mathbf{N})\) 在 X，第 68 页，注记 2）以及 X，第 87 页，注记 2）中引入。

定理 1. --- a) 若分解 R 或 S 中有一个是平坦的，则 \(\psi(S, R)\) 是 \(k\) -分次模的同构。

\begin{enumerate}
\def\labelenumi{\alph{enumi})}
\setcounter{enumi}{1}
\item
  若 R 是投射的，或 E 是内射的，则 \(\varphi(\mathbf{R},\mathbf{E})\) 是 \(k\) -分次模的同构。
\item
  例如假设 R 是平坦的，并如上选择 \(\alpha\) 与 \(\beta\) 。同态 \(\beta \otimes \alpha\) 是以下态射的复合
\end{enumerate}

\[
\mathrm{L} (\mathrm{P}) \otimes_ {\mathrm{A}} \mathrm{L} (\mathrm{M}) \xrightarrow {l _ {\mathrm{L} (\mathrm{P})} \otimes \alpha} \mathrm{L} (\mathrm{P}) \otimes_ {\mathrm{A}} \mathrm{R} \xrightarrow {\beta \otimes 1 _ {\mathrm{R}}} \mathrm{S} \otimes_ {\mathrm{A}} \mathrm{R}.
\]

由于 L(P)（相应地，R）是平坦的，且 \(\alpha\) （相应地， \(\beta\) ）是同调同构，故 \(1_{L(P)} \otimes \alpha\) （相应地， \(\beta \otimes 1_R\) ）也是同调同构（X，第 67 页，命题 4）。因此 \(\beta \otimes \alpha\) 是同调同构，且 \(\psi(S, R) = H(\beta \otimes \alpha)\) 是双射。

\begin{enumerate}
\def\labelenumi{\alph{enumi})}
\setcounter{enumi}{1}
\tightlist
\item
  类似地论证，使用 X，第 86 页，命题 4。
\end{enumerate}

推论. --- 若 R 是 M 的平坦分解，则同态

\[
\psi (\mathrm{P}, \mathrm{R}): \operatorname{Tor} ^ {\mathbf {A}} (\mathrm{P}, \mathrm{M}) \rightarrow \mathrm{H} (\mathrm{P} \otimes_ {\mathbf {A}} \mathrm{R})
\]

是双射。若 R 是 M 的投射分解，则同态

\[
\varphi (\mathrm{R}, \mathrm{N}): \mathrm{H} \left(\operatorname{Homgr} _ {\mathrm{A}} (\mathrm{R}, \mathrm{N})\right)\rightarrow \operatorname{Ext} _ {\mathrm{A}} (\mathrm{M}, \mathrm{N})
\]

是双射。若 E 是 N 的内射分解，则同态

\[
\varphi (\mathbf {M}, \mathbf {E}): \mathrm{H} \left(\operatorname{Homgr} _ {\mathbf {A}} (\mathbf {M}, \mathbf {E})\right)\rightarrow \operatorname{Ext} _ {\mathbf {A}} (\mathbf {M}, \mathbf {N})
\]

是双射。

注记. --- k-模的图表

\[
\begin{array}{c} \operatorname{Tor} ^ {\mathbf {A}} (\mathrm{P}, \mathrm{M}) \xrightarrow {\psi (\mathrm{S} , \mathrm{R})} \mathrm{H} (\mathrm{S} \otimes_ {\mathbf {A}} \mathrm{R}) \\ \sigma_ {\mathrm{P}, \mathrm{M}} \Biggl \downarrow \\ \operatorname{Tor} ^ {\mathbf {A} ^ {\circ}} (\mathrm{M} ^ {\circ}, \mathrm{P} ^ {\circ}) \xrightarrow {\psi (\mathrm{R} ^ {\circ} , \mathrm{S} ^ {\circ})} \mathrm{H} (\mathrm{R} ^ {\circ} \otimes_ {\mathbf {A} ^ {\circ}} \mathrm{S} ^ {\circ}), \end{array}
\]

其中 \(\sigma_{\mathrm{P,M}}\) 与 \(\sigma(\mathrm{S},\mathrm{R})\) 是交换同构（X，第 71 页和第 63 页），该图表是交换的：实际上，它是通过取同调，由复形的交换图

\[
\begin{array}{c c} \mathrm{L} (\mathrm{P}) \otimes_ {\mathrm{A}} \mathrm{L} (\mathrm{M}) & \xrightarrow {\beta \otimes \alpha} \mathrm{S} \otimes_ {\mathrm{A}} \mathrm{R} \\ \sigma (\mathrm{L} (\mathrm{P}), \mathrm{L} (\mathrm{M})) \Big \downarrow & \Big \downarrow \sigma (\mathrm{S}, \mathrm{R}) \\ \mathrm{L} (\mathrm{M} ^ {\circ}) \otimes_ {\mathrm{A} ^ {\circ}} \mathrm{L} (\mathrm{P} ^ {\circ}) & \xrightarrow {\alpha \otimes \beta} \mathrm{R} ^ {\circ} \otimes_ {\mathrm{A} ^ {\circ}} \mathrm{S} ^ {\circ} \end{array}
\]

（“态射 \(\psi\) 与交换同构相容”）。

类似地，设 \(a_{1}: \mathbb{R}_{1} \to \mathbb{M}\) , \(b_{1}: \mathbb{S}_{1} \to \mathbb{P}\) , \(c_{1}: \mathbb{N} \to \mathbb{E}_{1}\) 为复形态射，其中 \(\mathbb{R}_{1}\) 与 \(\mathbb{S}_{1}\) 是投射的且在右侧为零，并且 \(\mathbb{E}_{1}\) 是内射的且在左侧为零。根据 X，第 49 页，命题 3 和 3 bis，存在复形态射

\[
\alpha_ {1}: \mathrm{R} _ {1} \rightarrow \mathrm{L} (\mathrm{M}), \quad \beta_ {1}: \mathrm{S} _ {1} \rightarrow \mathrm{L} (\mathrm{P}), \quad \gamma_ {1}: \mathrm{I} (\mathrm{N}) \rightarrow \mathrm{E} _ {1}
\]

使得 \(p_{\mathrm{M}} \circ \alpha_{1} = a_{1}, p_{\mathrm{P}} \circ \beta_{1} = b_{1}, \gamma_{1} \circ e_{\mathrm{N}} = c_{1}\) 由此得到复形态射

\[
\beta_ {1} \otimes \alpha_ {1}: S _ {1} \otimes_ {A} R _ {1} \rightarrow L (P) \otimes_ {A} L (M),
\]

\[
\operatorname{Homgr} _ {\mathbf {A}} \left(\alpha_ {1}, \gamma_ {1}\right): \operatorname{Homgr} _ {\mathbf {A}} \left(\mathrm{L} (\mathrm{M}), \mathrm{I} (\mathrm{N})\right)\rightarrow \operatorname{Homgr} _ {\mathbf {A}} \left(\mathrm{R} _ {1}, \mathrm{E} _ {1}\right),
\]

并且通过取同调，得到次数为 0 的 k-线性分次映射：

\[
\begin{array}{l} \psi^ {\prime} (\mathbf {S} _ {1}, \mathbf {R} _ {1}): \mathrm{H} (\mathbf {S} _ {1} \otimes_ {\mathbf {A}} \mathbf {R} _ {1}) \to \operatorname{Tor} ^ {\mathbf {A}} (\mathbf {P}, \mathbf {M}) \\ \varphi^ {\prime} (\mathbf {R} _ {1}, \mathbf {E} _ {1}): \operatorname{Ext} _ {\mathbf {A}} (\mathbf {M}, \mathbf {N}) \to \mathrm{H} (\operatorname{Homgr} _ {\mathbf {A}} (\mathbf {R} _ {1}, \mathbf {E} _ {1})) \end{array}
\]

如上所述，可验证这些映射不依赖于 \(\alpha_{1}\) , \(\beta_{1}\) , \(\gamma_{1}\) 的选择。

命题 1. --- 若 \(a_{1}\) , \(b_{1}\) , \(c_{1}\) 是同调同构，则 \(\psi^{\prime}(\mathbf{S}_{1},\mathbf{R}_{1})\) 与 \(\varphi^{\prime}(\mathbf{R}_{1},\mathbf{E}_{1})\) 分别是双射 \(\psi(\mathbf{S}_{1},\mathbf{R}_{1})\) 与 \(\varphi(\mathbf{R}_{1},\mathbf{E}_{1})\) 的逆双射。事实上， \(f = (\beta \otimes \alpha) \circ (\beta_{1} \otimes \alpha_{1})\) 是复形 \(\mathbf{S}_{1} \otimes_{\mathbf{A}} \mathbf{R}_{1}\) 到自身的态射，且有 \((b_{1} \circ \alpha_{1}) \circ f = f\) 。根据 X，第 49 页，命题 3 和 3 bis， \(f\) 是同伦等价，因此 H(f) = 1 且

\[
\psi \left(\mathrm{S} _ {1}, \mathrm{R} _ {1}\right) \circ \psi^ {\prime} \left(\mathrm{S} _ {1}, \mathrm{R} _ {1}\right) = \mathrm{H} (\beta \otimes \alpha) \circ \mathrm{H} \left(\beta_ {1} \otimes \alpha_ {1}\right) = \mathrm{H} (f) = 1;
\]

同理 \(\psi'(S_1, R_1) \circ \psi(S_1, R_1) = 1\) 。对映射 \(\varphi\) 与 \(\varphi'\) 可类似论证。

例。--- 1) 设 \(a\) 是 A 的一个元素，使得映射 \(\varphi : x \mapsto xa\) 从 A 到自身是单射（ \(a\) 不是右零因子）。利用分解

\[
0 \rightarrow \mathrm{A} _ {s} \xrightarrow {\varphi} \mathrm{A} _ {s} \rightarrow \mathrm{A} / \mathrm{A} a \rightarrow 0
\]

我们看到，对于每个右 A-模 M，有

\[
\operatorname{Tor} _ {i} ^ {\mathbf {A}} (\mathbf {M}, \mathbf {A} / \mathbf {A} a) = 0 \quad \text { for } \quad i > 1
\]

并且该 \(k\) -模 \(\operatorname{Tor}_1^{\mathbf{A}}(\mathbf{M}, \mathbf{A} / \mathbf{A}a)\) 同构于 \(\operatorname{Ker}(a_{\mathbf{M}})\) .

类似地，对于每个左 A-模 M，有

\[
\operatorname{Ext} _ {\mathbf {A}} ^ {i} (\mathbf {A} / \mathbf {A} a, \mathbf {M}) = 0 \quad \text { for } \quad i > 1
\]

且 \(k\) -模 \(\operatorname{Ext}_{\mathbf{A}}^{1}(\mathbf{A}/\mathbf{A}a, \mathbf{M})\) 同构于 \(\mathbf{M}/a\mathbf{M}\) .

\begin{enumerate}
\def\labelenumi{\arabic{enumi})}
\setcounter{enumi}{1}
\tightlist
\item
  设 A 是一个整环；令 K 为 A 的分式域，M 为一个 A-模。利用平坦分解
\end{enumerate}

\[
0 \rightarrow \mathrm{A} \rightarrow \mathrm{K} \rightarrow \mathrm{K} / \mathrm{A} \rightarrow 0
\]

（X，第 9 页，例 5），我们看到 \(\operatorname{Tor}_{i}^{\mathbf{A}}(\mathbf{K}/\mathbf{A}, \mathbf{M}) = 0\) 对 \(i > 1\) 成立；此外，考虑到 II，第 116 页，命题 26, (ii)，A-模 \(\operatorname{Tor}_{1}^{\mathbf{A}}(\mathbf{K}/\mathbf{A}, \mathbf{M})\) 同构于 M 的挠子模。

\begin{enumerate}
\def\labelenumi{\arabic{enumi})}
\setcounter{enumi}{2}
\tightlist
\item
  设 A 是一个诺特局部环，用 m 表示其极大理想，并令 \(\kappa = \mathbf{A} / \mathfrak{m}\) 。设 M 为有限生成的 A-模，P 为 M 的一个极小投射分解（X，第 54 页）。对每个 \(n \geqslant 0\) ， \(\kappa\) -向量空间 \(\operatorname{Tor}_n^{\mathbf{A}}(\kappa, \mathbf{M})\) 与 \(\operatorname{Ext}_{\mathbf{A}}^{n}(\mathbf{M}, \kappa)\) 都是有限维的，其维数等于自由 A-模 \(\mathbf{P}_n\) 的秩；事实上，复形 \(\kappa \otimes_{\mathbf{A}} \mathbf{P}\) 与 \(\operatorname{Homgr}_{\mathbf{A}}(\mathbf{P}, \kappa)\) 的微分均为零。
\end{enumerate}

